\section*{Capítulo \ref{ch:g12:climate-change} --- A mudança do clima: observações, causas e projeções}

\begin{solution}{exo:g12:climate-change:1}
A variação que ela causa no fluxo líquido de energia que entra no sistema
climático no topo da atmosfera, antes que a temperatura da superfície
responda; medido em \unit{W/m^2}. Positivo: o planeta ganha energia e se
aquece; negativo: ele perde energia e se resfria.
\end{solution}

\begin{solution}{exo:g12:climate-change:2}
Máxima em maio, mínima em setembro. Na primavera e no verão, a vegetação do
Hemisfério Norte, onde cresce a maior parte das plantas terrestres, absorve
pela fotossíntese mais gás carbônico do que as plantas e os solos liberam;
no outono e no inverno, a decomposição e a respiração dominam e a
concentração volta a subir.
\end{solution}

\begin{solution}{exo:g12:climate-change:3}
\emph{Muito provável}: uma probabilidade de 90 a 100\,\%; \emph{provável}:
de 66 a 100\,\%. $[0.95 \text{ a } 1.20]$ é uma faixa muito provável: 90\,\%
de chance de que o valor verdadeiro esteja entre 0.95 e 1.20.
\end{solution}

\begin{solution}{exo:g12:climate-change:4}
91\,\%. A água se dilata quando se aquece: um oceano mais quente ocupa mais
volume, então a sua superfície sobe mesmo que nenhuma água seja
acrescentada (o gelo continental que derrete acrescenta água além disso).
\end{solution}

\begin{solution}{exo:g12:climate-change:5}
O vapor de água: um ar mais quente contém mais vapor de água, um gás de
efeito estufa, o que reforça o aquecimento. A neve e o gelo: o derretimento
descobre superfícies mais escuras, que absorvem mais luz do Sol. A resposta
negativa: um planeta mais quente irradia mais infravermelho para o espaço,
o que o traz de volta ao equilíbrio.
\end{solution}

\begin{solution}{exo:g12:climate-change:6}
$427.35 - 315.98 = \qty{111.37}{ppm}$; $111.37/315.98 \approx 0.352$, um
aumento de cerca de 35\,\%. Em 66 anos: $111.37/66 \approx
\qty{1.7}{ppm}$ por ano, em média.
\end{solution}

\begin{solution}{exo:g12:climate-change:7}
$F = 3.93 \times \ln(427.35/278)/\ln 2 = 3.93 \times 0.4300/0.6931 \approx
\qty{2.44}{W/m^2}$.
\end{solution}

\begin{solution}{exo:g12:climate-change:8}
$3.7 \times 30 = \qty{111}{mm}$, cerca de \qty{11}{cm}. $3.7/1.3 \approx 2.8$:
quase três vezes mais rápida.
\end{solution}

\begin{solution}{exo:g12:climate-change:9}
$\Delta T = 2.72/1.16 \approx \qty{2.3}{\degreeCelsius}$. O clima não está
em equilíbrio: o oceano, que absorve cerca de 90\,\% da energia extra, se
aquece devagar, então a superfície ainda não alcançou o \omterm{def:g12:climate-change:forcing}{forçamento} (o
sistema ainda ganha cerca de \qty{0.8}{W/m^2}).
\end{solution}

\begin{solution}{exo:g12:climate-change:10}
O Sol e os vulcões explicam no máximo \qty{0.1}{\degreeCelsius} de um
aquecimento de \qty{1.06}{\degreeCelsius}: eles não podem ser a sua causa.
A outra causa natural é a variabilidade interna do clima (as suas
flutuações naturais, como as trocas de calor entre o oceano e a atmosfera);
a avaliação estima a sua contribuição em $-0.2$ a $+\qty{0.2}{\degreeCelsius}$,
também pequena demais.
\end{solution}

\begin{solution}{exo:g12:climate-change:11}
Muito provavelmente abaixo de \qty{2}{\degreeCelsius}: só o SSP1-1.9 (faixa
muito provável de 1.0 a 1.8). Muito provavelmente acima: o SSP2-4.5 (2.1 a
3.5), o SSP3-7.0 (2.8 a 4.6) e o SSP5-8.5 (3.3 a 5.7). Para o SSP1-2.6 (1.3
a 2.4), nenhum dos dois é muito provável.
\end{solution}

\begin{solution}{exo:g12:climate-change:12}
O \omterm{def:g12:climate-change:forcing}{forçamento} dos aerossóis é $-0.22 - 0.84 = \qty{-1.06}{W/m^2}$; sem ele, o
total passa a $2.72 + 1.06 = \qty{3.78}{W/m^2}$. O aquecimento de equilíbrio
aumenta $1.06/1.16 \approx \qty{0.9}{\degreeCelsius}$.
\end{solution}

\begin{solution}{exo:g12:climate-change:13}
Como $\ln(2C) - \ln C = \ln 2$ qualquer que seja $C$,
\[
  F(556) - F(278) = \frac{3.93\,\ln 2}{\ln 2} = \qty{3.93}{W/m^2}, \qquad
  F(1112) - F(556) = \frac{3.93\,\ln 2}{\ln 2} = \qty{3.93}{W/m^2}:
\]
o mesmo \omterm{def:g12:climate-change:forcing}{forçamento} para cada duplicação. Total a \qty{1112}{ppm}: $2 \times 3.93 =
\qty{7.86}{W/m^2}$.
\end{solution}

\begin{solution}{exo:g12:climate-change:14}
$\dfrac{\dd F}{\dd C} = \dfrac{F_{2\times}}{C \ln 2}$. A \qty{280}{ppm}:
$3.93/(280 \times 0.6931) \approx \qty{0.0202}{W/m^2}$ por ppm; a
\qty{420}{ppm}: $\approx \qty{0.0135}{W/m^2}$ por ppm. Cada ppm extra
acrescenta menos \omterm{def:g12:climate-change:forcing}{forçamento} à medida que a concentração cresce (1.5 vez
menos a 420 que a 280), mas a concentração continua subindo, então o
\omterm{def:g12:climate-change:forcing}{forçamento} continua crescendo.
\end{solution}

\begin{solution}{exo:g12:climate-change:15}
$3.93/3.22 \approx \qty{1.2}{\degreeCelsius}$ sem \omterm{def:g12:climate-change:feedback}{retroalimentações};
$3.93/1.16 \approx \qty{3.4}{\degreeCelsius}$ com elas. Fator de
amplificação: $3.22/1.16 \approx 2.8$.
\end{solution}

\begin{solution}{pb:g12:climate-change:1}
\textbf{1.} $3.37 + 3.51 = \qty{6.88}{ppm}$, cerca de \qty{7}{ppm}.

\textbf{2.} Da primavera ao fim do verão, as plantas do Hemisfério Norte
retiram gás carbônico do ar pela fotossíntese mais depressa do que a
respiração e a decomposição o devolvem: a concentração cai até setembro. No
outono e no inverno, o balanço se inverte, e a concentração volta a subir
até o seu máximo de maio.

\textbf{3.} O Hemisfério Sul tem muito menos terras, e portanto muito menos
vegetação, que o Norte; a maior parte da sua superfície é oceano.

\textbf{4.} No alto, ele fica acima da camada de ar poluído perto do solo;
longe dos continentes e das cidades, o ar que ele amostra é bem misturado e
representa a atmosfera inteira, e não uma fonte local.

\textbf{5.} $(325.68 - 316.91)/10 \approx \qty{0.88}{ppm}$ por ano.

\textbf{6.} $(427.35 - 401.01)/10 \approx \qty{2.63}{ppm}$ por ano.

\textbf{7.} $2.634/0.877 \approx 3.0$: a taxa de crescimento triplicou.

\textbf{8.} $(427.35 - 278)/278 \approx 0.537$: cerca de 54\,\%.

\textbf{9.} $427.35 - 300 \approx \qty{127}{ppm}$.

\textbf{10.} $F = 3.93 \times \ln(427.35/278)/\ln 2 \approx
\qty{2.44}{W/m^2}$.

\textbf{11.} $F(316.91) = 3.93 \times \ln(316.91/278)/\ln 2 \approx
\qty{0.74}{W/m^2}$; aumento de $2.44 - 0.74 \approx \qty{1.70}{W/m^2}$.

\textbf{12.} Um \unit{W/m^2} a mais multiplica a concentração por
$2^{1/3.93} \approx 1.193$: $427.35 \times 1.193 \approx \qty{510}{ppm}$.

\textbf{13.} $(509.8 - 427.35)/2.634 \approx 31$ anos.

\textbf{14.} $2 \times 278 = \qty{556}{ppm}$; \omterm{def:g12:climate-change:forcing}{forçamento} $F_{2\times} =
\qty{3.93}{W/m^2}$.

\textbf{15.} Pela definição da sensibilidade, \qty{3}{\degreeCelsius}
(provavelmente de 2.5 a \qty{4}{\degreeCelsius}).

\textbf{16.} $\lambda = 3.93/3 \approx \qty{1.31}{W.m^{-2}.\degreeCelsius^{-1}}$.
A soma das \omterm{def:g12:climate-change:feedback}{retroalimentações}, \qty{1.16}{W.m^{-2}.\degreeCelsius^{-1}}, dá
$3.93/1.16 \approx \qty{3.4}{\degreeCelsius}$, um pouco mais; a melhor
estimativa avaliada combina essa física com o aquecimento observado e com os
climas do passado, e os dois valores ficam na faixa provável.

\textbf{17.} O oceano armazena a maior parte da energia extra e se aquece
devagar: a superfície só atinge o seu aquecimento de equilíbrio depois de
décadas a séculos.

\textbf{18.} Ele supõe uma história das emissões futuras (o crescimento dos
últimos dez anos mantido para sempre): isso é um \omterm{def:g12:climate-change:scenario}{cenário}. Uma \omterm{def:g12:climate-change:scenario}{projeção}
calcula a resposta do clima e do ciclo do carbono a um \omterm{def:g12:climate-change:scenario}{cenário de emissões}
com modelos físicos, e dá faixas e probabilidades.

\textbf{19.} $(556 - 427.35)/2.634 \approx 48.8$ anos depois de 2025: por
volta de \textbf{2074}.
\end{solution}
